\chapter{Agents and the Scientific Method}
\label{sec:agents}
% ===================================================================

The agent defined below is thinner than the one
Chapter~\ref{ch:sci} built, deliberately so.
That one had a purse, a stratified distance on its hypothesis space, a
cost for every experiment, and a demonstrable tendency to prefer false
theories that were cheap to true ones that were not.
This one has a hypothesis language, a world model, and a revision rule,
and it is not going to starve.

The thinness is the point.
The claim of this part is that a population of agents fails to reach the
true ontology \emph{even when the agents are ideal} --- even when the
scientific method is executed perfectly and convergence is guaranteed in
the limit.
Introducing a budget here would make the conclusion easy and
uninteresting, since a budgeted agent obviously cannot afford everything.
What is interesting is that the obstruction survives the removal of the
budget, and reappears as a fact about communication geometry rather than
about anyone's purse.

\begin{remark}[Which agent is doing the work]
\label{rmk:two-agents}
The two agents are not rivals and neither supersedes the other.
The mortal scientist of Part~\ref{part:ecology} is what an agent
\emph{is} in this framework; the idealised agent below is a limiting case
of it, obtained by letting the purse grow without bound.
Where a later chapter needs an agent that can be misled, it is this one;
where it needs an agent that can die, it is that one.
Chapter~\ref{sec:false-ontology} brings the budget back at the end, and
the reintroduction is what turns a claim about limits into a claim about
what any actual scientist will infer.
\end{remark}

\bigskip

\section{Agents as Terms}

\begin{definition}[Agent]
An \emph{agent} in a GSLT $\GSLT$ is a term $A \in \terms(\GSLT)$
equipped with:
\begin{enumerate}[label=(\roman*)]
  \item a \emph{hypothesis language}: the context-decorated
        HML $\HML(\ctx)$ derived from $\GSLT$;
  \item a \emph{world model}: a subset $\mathcal{W} \subseteq
        \terms(\GSLT)/{\bisim}$ representing the agent's current
        best approximation to the bisimulation quotient;
  \item an \emph{experimental budget}: a vectorial account
        $\vect{A}$ as defined in Part~I.
\end{enumerate}
\end{definition}

\begin{remark}
The hypothesis language and experimental budget are structures
already present in Part~I.
The world model is new: it is the agent's \emph{internal} representation
of its environment, constructed and refined through experiment.
\end{remark}

\section{The Scientific Method as a Fixed-Point Iteration}

We model the scientific method as the following iterative procedure,
which the agent runs within its experimental budget:

\begin{enumerate}[label=\textbf{Step \arabic*.}, leftmargin=*, itemsep=4pt]
  \item \textbf{Hypothesize.} Construct a formula
        $\phi \in \HML(\ctx)$ that makes a testable claim about
        some term $T$ in the environment.
        A formula $\phi$ is \emph{testable} if there exists a
        context $K$ such that $\phi = \langle K \rangle \psi$
        for some $\psi$ --- i.e., the formula has a witness context.

  \item \textbf{Design experiment.} Extract the witness context $K$
        from $\phi$.
        By the Milner--Sewell--Leifer construction, minimal contexts
        and HML formulae are in correspondence; the context $K$ is
        the \emph{experimental apparatus} that tests $\phi$.

  \item \textbf{Run experiment.} Form $K[T]$ and observe whether it
        reduces to a term satisfying $\psi$.
        Debit the cost $\vect{c}$ of the interaction from the account
        $\vect{A}$.

  \item \textbf{Update world model.} If the experiment confirms $\phi$,
        refine the equivalence class of $T$ in $\mathcal{W}$ to include
        the constraint $T \sat \phi$.
        If it refutes $\phi$, add $T \sat \neg\phi$.

  \item \textbf{Repeat} until the budget is exhausted or the world
        model stabilizes.
\end{enumerate}

\begin{proposition}[Convergence in the Limit]
If the experimental budget is unlimited and the GSLT has decidable
bisimilarity, the scientific method converges: the world model
$\mathcal{W}$ stabilizes at the bisimulation quotient
$\terms(\GSLT)/{\bisim}$.
\end{proposition}

\begin{remark}
The convergence proposition is the idealized case.
For GSLTs with undecidable bisimilarity (e.g., the full
$\pi$-calculus), convergence is not guaranteed even with
unlimited budget.
This is the computational analogue of undecidable physical questions ---
there are aspects of the world the agent can never experimentally resolve,
regardless of resources.
The present paper is concerned with a different and more fundamental
obstruction to convergence: the communication geometry of a massive
population, which limits the budget effectively to zero for the
vast majority of the population.
\end{remark}

% ===================================================================
